\documentclass{stex}

\libusepackage{naproche}
\libinput{preamble}

\begin{document}
\begin{smodule}{ordering.ftl}

\importmodule[libraries/arithmetics]{natural-numbers.ftl}

\symdef{Less}[name=less]{\;<\;}
\symdef{NotLess}{\;\not<\;}
\symdef{Greater}[name=greater]{\;>\;}
\symdef{NotGreater}{\;\not>\;}
\symdef{Leq}{\;\leq\;}
\symdef{NotLeq}{\;\not\leq\;}
\symdef{Geq}{\;\geq\;}
\symdef{NotGeq}{\;\not\geq\;}
\symdef{NatLess}[args=1]{\comp{\mathbb{N}}_{\comp<#1}}
\symdef{NatGreater}[args=1]{\comp[\mathbb{N}]_{\comp>#1}}
\symdef{NatLeq}[args=1]{\comp{\mathbb{N}}_{\comp\leq#1}}
\symdef{NatGeq}[args=1]{\comp{\mathbb{N}}_{\comp\geq#1}}

\symdecl*{positive}
\symdecl*{predecessor}
\symdecl*{successor}
\symdecl*{preservation of ordering under right-addition}
\symdecl*{wellfoundedness of ordering}
\symdecl*{induction II}
\symdecl*{induction III}

\begin{definition*}[forthel,for={less,NotLess,greater,NotGreater}]
  Let $n, m$ be natural numbers.
  $\emph{n \Less m}$ iff there exists a nonzero natural number $k$ such that $m \Eq n \Plus k$.
  Let $n$ is \emph{less than $m$} stand for $n \Less m$.
  Let $\emph{n \Greater m}$ stand for $m \Less n$.
  Let $n$ is \emph{greater than $m$} stand for $n \Greater m$.
  Let $\emph{n \NotLess m}$ stand for $n$ is not less than $m$.
  Let $\emph{n \NotGreater m}$ stand for $n$ is not greater than $m$.
\end{definition*}

\begin{definition*}[forthel,for=NatLess]
  Let $n$ be a natural number.
  $\emph{\NatLess{n}} \DefEq \SClass{k}{\Naturals}{k \Less n}$.
\end{definition*}

\begin{definition*}[forthel,for=NatGreater]
  Let $n$ be a natural number.
  $\emph{\NatGreater{n}} \DefEq \SClass{k}{\Naturals}{k \Greater n}$.
\end{definition*}

\begin{definition*}[forthel,for=positive]
  Let $n$ be a natural number.
  $n$ is \emph{positive} iff $n \Greater \Zero$.
\end{definition*}

\begin{definition*}[forthel,for={Leq,Geq,NotLeq,NotGeq}]
  Let $n, m$ be natural numbers.
  $\emph{n \Leq m}$ iff there exists a natural number $k$ such that $m \Eq n \Plus k$.
  Let $n$ is \emph{less than or equal to $m$} stand for $n \Leq m$.
  Let $\emph{n \Geq m}$ stand for $m \Leq n$.
  Let $n$ is \emph{greater than or equal to $m$} stand for $n \Geq m$.
  Let $\emph{n \NotLeq m}$ stand for $n$ is not less than or equal to $m$.
  Let $\emph{n \NotGeq m}$ stand for $n$ is not greater than or equal to $m$.
\end{definition*}

\begin{definition*}[forthel,for=NatLeq]
  Let $n$ be a natural number.
  $\emph{\NatLeq{n}} \DefEq \SClass{k}{\Naturals}{k \Leq n}$.
\end{definition*}

\begin{definition*}[forthel,for=NatGeq]
  Let $n$ be a natural number.
  $\emph{\NatGeq{n}} \DefEq \SClass{k}{\Naturals}{k \Geq n }$.
\end{definition*}

\begin{proposition*}[forthel]
  Let $n, m$ be natural numbers.
  $n \Leq m$ iff $n \Less m$ or $n \Eq m$.
\end{proposition*}
\begin{proof}[forthel]
  \begin{case}{$n \Leq m$.}
    Take a natural number $k$ such that $m \Eq n \Plus k$.
    If $k \Eq \Zero$ then $n \Eq m$. If $k \NotEq \Zero$ then $n \Less m$.
  \end{case}

  \begin{case}{$n \Less m$ or $n \Eq m$.}
    If $n \Less m$ then there is a positive natural number $k$ such that $m \Eq n \Plus k$.
    If $n \Eq m$ then $m \Eq n \Plus \Zero$.
    Thus if $n \Less m$ then there is a natural number $k$ such that $m \Eq n \Plus k$.
  \end{case}
\end{proof}

\begin{definition*}[forthel,for=predecessor]
  Let $n$ be a natural number.
  A \emph{predecessor of $n$} is a natural number that is less than $n$.
\end{definition*}

\begin{definition*}[forthel,for=successor]
  Let $n$ be a natural number.
  A \emph{successor of $n$} is a natural number that is greater than $n$.
\end{definition*}

\begin{proposition*}[forthel]
  Let $n$ be a natural number.
  Then $n$ is positive iff $n$ is nonzero.
\end{proposition*}
\begin{proof}[forthel]
  \begin{case}{$n$ is positive.}
    Take a positive natural number $k$ such that $n \Eq \Zero \Plus k \Eq k$.
    Then we have $n \NotEq \Zero$.
  \end{case}

  \begin{case}{$n$ is nonzero.}
    Take a natural number $k$ such that $n \Eq k \Plus \One$.
    Then $n \Eq \Zero \Plus (k \Plus \One)$.
    $k \Plus \One$ is positive.
    Hence $\Zero \Less n$.
  \end{case}
\end{proof}

\begin{proposition*}[forthel]
  Let $n$ be a natural number.
  Then \[ n \NotLess n. \]
\end{proposition*}
\begin{proof}[forthel]
  Assume the contrary.
  Then we can take a positive natural number $k$ such that $n \Eq n \Plus k$.
  Then we have $\Zero \Eq k$.
  Contradiction.
\end{proof}

\begin{proposition*}[forthel]
  Let $n, m$ be natural numbers.
  Then \[ n \Less m \Implies n \NotEq m. \]
\end{proposition*}
\begin{proof}[forthel]
  Assume $n \Less m$.
  Take a positive natural number $k$ such that $m \Eq n \Plus k$.
  If $n \Eq m$ then $k \Eq \Zero$.
  Hence $n \NotEq m$.
\end{proof}

\begin{proposition*}[forthel]
  Let $n, m$ be natural numbers.
  If $n \Leq m$ and $m \Leq n$ then $n \Eq m$.
\end{proposition*}
\begin{proof}[forthel]
  Assume $n \Leq m$ and $m \Leq n$.
  Take natural numbers $k, l$ such that $m \Eq n \Plus k$ and $n \Eq m \Plus l$.
  Then $m
    \Eq n \Plus k
    \Eq (m \Plus l) \Plus k
    \Eq m \Plus (l \Plus k)$.
  Hence $l \Plus k \Eq \Zero$.
  Thus $l \Eq \Zero \Eq k$.
  Indeed if $l \NotEq \Zero$ or $k \NotEq \Zero$ then $l \Plus k$ is the direct successor of
  some natural number.
  Therefore $m \Eq n$.
\end{proof}

\begin{proposition*}[forthel]
  Let $n, m, k$ be natural numbers.
  If $n \Less m \Less k$ then $n \Less k$.
\end{proposition*}
\begin{proof}[forthel]
  Assume $n \Less m \Less k$.
  Take a positive natural number $a$ such that $m \Eq n \Plus a$.
  Take a positive natural number $b$ such that $k \Eq m \Plus b$.
  Then $k
    \Eq m \Plus b
    \Eq (n \Plus a) \Plus b
    \Eq n \Plus (a \Plus b)$.
  $a \Plus b$ is positive.
  Hence $n \Less k$.
\end{proof}

\begin{proposition*}[forthel]
  Let $n, m, k$ be natural numbers.
  If $n \Leq m \Leq k$ then $n \Leq k$.
\end{proposition*}
\begin{proof}[forthel]
  Assume $n \Leq m \Leq k$.
  \begin{case}{$n \Eq m \Eq k$.} \end{case}
  \begin{case}{$n \Eq m \Less k$.} \end{case}
  \begin{case}{$n \Less m \Eq k$.} \end{case}
  \begin{case}{$n \Less m \Less k$.} \end{case}
\end{proof}

\begin{proposition*}[forthel]
  Let $n, m, k$ be natural numbers.
  If $n \Leq m \Less k$ then $n \Less k$.
\end{proposition*}
\begin{proof}[forthel]
  Assume $n \Leq m \Less k$.
  If $n \Eq m$ then $n \Less k$.
  If $n \Less m$ then $n \Less k$.
\end{proof}

\begin{proposition*}[forthel]
  Let $n, m, k$ be natural numbers.
  If $n \Less m \Leq k$ then $n \Less k$.
\end{proposition*}
\begin{proof}[forthel]
  Assume $n \Less m \Leq k$.
  If $m \Eq k$ then $n \Less k$.
  If $m \Less k$ then $n \Less k$.
\end{proof}

\begin{proposition*}[forthel]
  Let $n, m$ be natural numbers.
  If $n \Less m$ then $n \Plus \One \Leq m$.
\end{proposition*}
\begin{proof}[forthel]
  Assume $n \Less m$.
  Take a positive natural number $k$ such that $m \Eq n \Plus k$.

  \begin{case}{$k \Eq \One$.}
    Then $m \Eq n \Plus \One$.
    Hence $n \Plus \One \Leq m$.
  \end{case}

  \begin{case}{$k \NotEq \One$.}
    Then we can take a natural number $l$ such that $k \Eq l \Plus \One$.
    Then $m
      \Eq n \Plus (l \Plus \One)
      \Eq (n \Plus l) \Plus \One
      \Eq (n \Plus \One) \Plus l$.
    $l$ is positive.
    Thus $n \Plus \One \Less m$.
  \end{case}
\end{proof}

\begin{proposition*}[forthel]
  Let $n, m$ be natural numbers.
  Then $n \Less m$ or $n \Eq m$ or $n \Greater m$.
\end{proposition*}
\begin{proof}[forthel]
  Define $\Phi \DefEq \SClass{m'}{\Naturals}{n \Less m'\text{ or }n \Eq m'\text{ or }n \Greater m'}$.

  (1) $\Phi$ contains $\Zero$.

  (2) For all $m' \Elem \Phi$ we have $m' \Plus \One \Elem \Phi$.
  \begin{proof}
    Let $m' \Elem \Phi$.

    \begin{case}{$n \Less m'$.} \end{case}

    \begin{case}{$n \Eq m'$.} \end{case}

    \begin{case}{$n \Greater m'$.}
      Take a positive natural number $k$ such that $n \Eq m' \Plus k$.

      \begin{case}{$k \Eq \One$.} \end{case}

      \begin{case}{$k \NotEq \One$.}
        Take a natural number $l$ such that $n \Eq (m' \Plus \One) \Plus l$.
        Hence $n \Greater m' \Plus \One$.
        Indeed $l$ is positive.
      \end{case}
    \end{case}
  \end{proof}

  Thus every natural number is contained in $\Phi$ (by \sr{induction I}{induction}).
  Therefore $n \Less m$ or $n \Eq m$ or $n \Greater m$.
\end{proof}

\begin{proposition*}[forthel]
  Let $n, m$ be natural numbers.
  Then $n \NotLess m$ iff $n \Geq m$.
\end{proposition*}
\begin{proof}[forthel]
  \begin{case}{$n \NotLess m$.}
    Then $n \Eq m$ or $n \Greater m$.
    Hence $n \Geq m$.
  \end{case}

  \begin{case}{$n \Geq m$.}
    Assume $n \Less m$.
    Then $n \Leq m$.
    Hence $n \Eq m$.
    Contradiction.
  \end{case}
\end{proof}

\begin{proposition*}[forthel]
  Let $n, m$ be natural numbers.
  If $n \Less m \Leq n \Plus \One$ then $m \Eq n \Plus \One$.
\end{proposition*}
\begin{proof}[forthel]
  Assume $n \Less m \Leq n \Plus \One$.
  Take a positive natural number $k$ such that $m \Eq n \Plus k$.
  Take a natural number $l$ such that $n \Plus \One \Eq m \Plus l$.
  Then $n \Plus \One
    \Eq m \Plus l
    \Eq (n \Plus k) \Plus l
    \Eq n \Plus (k \Plus l)$.
  Hence $k \Plus l \Eq \One$.

  We have $l \Eq \Zero$.
  \begin{proof}
    Assume the contrary.
    Then $k,l \Greater \Zero$.

    \begin{case}{$k,l \Eq \One$.}
      Then $k \Plus l
        \Eq \Two
        \NotEq \One$.
      Contradiction.
    \end{case}

    \begin{case}{$k \Eq \One and l \NotEq \One$.}
      Then $l \Greater \One$.
      Hence $k \Plus l
        \Greater \One \Plus l
        \Greater \One$.
      Contradiction.
    \end{case}

    \begin{case}{$k \NotEq \One and l \Eq \One$.}
      Then $k \Greater \One$.
      Hence $k \Plus l
        \Greater k \Plus \One
        \Greater \One$.
      Contradiction.
    \end{case}

    \begin{case}{$k, l \NotEq \One$.}
      Take natural numbers $a, b$ such that $k \Eq a \Plus \One$ and $l \Eq b \Plus \One$.
      Indeed $k, l \NotEq \Zero$.
      Hence $k \Eq a \Plus \One$ and $l \Eq b \Plus \One$.
      Thus $k, l \Greater \One$. Indeed $a, b$ are positive.
    \end{case}
  \end{proof}

  Then we have $n \Plus \One
    \Eq m \Plus l
    \Eq m \Plus \Zero
    \Eq m$.
\end{proof}

\begin{proposition*}[forthel]
  Let $n, m$ be natural numbers.
  If $n \Leq m \Less n \Plus \One$ then $n \Eq m$.
\end{proposition*}
\begin{proof}[forthel]
  Assume $n \Leq m \Less n \Plus \One$.

  \begin{case}{$n \Eq m$.} \end{case}

  \begin{case}{$n \Less m$.}
    Then $n \Less m \Leq n \Plus \One$.
    Hence $n \Eq m$.
  \end{case}
\end{proof}

\begin{corollary*}[forthel]
  Let $n$ be a natural number.
  There is no natural number $m$ such that $n \Less m \Less n \Plus \One$.
\end{corollary*}
\begin{proof}[forthel]
  Assume the contrary.
  Take a natural number $m$ such that $n \Less m \Less n \Plus \One$.
  Then $n \Less m \Leq n \Plus \One$ and $n \Leq m \Less n \Plus \One$.
  Hence $m \Eq n \Plus \One$ and $m \Eq n$.
  Hence $n \Eq n \Plus \One$.
  Contradiction.
\end{proof}

\begin{proposition*}[forthel]
  Let $n$ be a natural number.
  Then $n \Plus \One \Geq \One$.
\end{proposition*}
\begin{proof}[forthel]
  \begin{case}{$n \Eq \Zero$.} \end{case}

  \begin{case}{$n \NotEq \Zero$.}
    Then $n \Greater \Zero$.
    Hence $n \Plus \One \Greater \Zero \Plus \One \Eq \One$.
  \end{case}
\end{proof}

\begin{proposition*}[forthel,name=preservation of ordering under right-addition]
  Let $n, m, k$ be natural numbers.
  Then $n \Less m$ iff $n \Plus k \Less m \Plus k$.
\end{proposition*}
\begin{proof}[forthel]
  \begin{case}{$n \Less m$.}
    Take a positive natural number $l$ such that $m \Eq n \Plus l$.
    Then $m \Plus k
      \Eq (n \Plus l) \Plus k
      \Eq (n \Plus k) \Plus l$.
    Hence $n \Plus k \Less m \Plus k$.
  \end{case}

  \begin{case}{$n \Plus k \Less m \Plus k$.}
    Take a positive natural number $l$ such that $m \Plus k \Eq (n \Plus k) \Plus l$.
    $((n \Plus k) \Plus l)
      \Eq n \Plus (k \Plus l)
      \Eq n \Plus (l \Plus k)
      \Eq (n \Plus l) \Plus k$.
    Hence $m \Plus k \Eq (n \Plus l) \Plus k$.
    Thus $m \Eq n \Plus l$ (by \sn{right-cancellability of addition}).
    Therefore $n \Less m$.
  \end{case}
\end{proof}

\begin{corollary*}[forthel]
  Let $n, m, k$ be natural numbers.
  Then $n \Less m$ iff $k \Plus n \Less k \Plus m$.
\end{corollary*}
\begin{proof}[forthel]
  We have $k \Plus n \Eq n \Plus k$ and $k \Plus m \Eq m \Plus k$.
  Hence $k \Plus n \Less k \Plus m$ iff $n \Plus k \Less m \Plus k$.
\end{proof}

\begin{corollary*}[forthel]
  Let $n, m, k$ be natural numbers.
  Then $n \Leq m$ iff $k \Plus n \Leq k \Plus m$.
\end{corollary*}

\begin{corollary*}[forthel]
  Let $n, m, k$ be natural numbers.
  Then $n \Leq m$ iff $n \Plus k \Leq m \Plus k$.
\end{corollary*}

\begin{proposition*}[forthel,name=wellfoundedness of ordering]
  Let $A$ be a nonempty subclass of $\Naturals$.
  Then there exists a $m \Elem A$ such that $m \Leq n$ for all $n \Elem A$.
\end{proposition*}
\begin{proof}[forthel]
  Assume the contrary.

  Let us show that for each $n \Elem A$ there exists a $m \Elem A$ such that $m \Less n$.
    Let $n \Elem A$.
    Assume that there exists no $m \Elem A$ such that $m \Less n$.
    Then $n \Leq m$ for all $m \Elem A$.
    Contradiction.
  End.

  (a) Define $\Phi \DefEq \SClass{n}{\Naturals}{n\text{ is less than any element of }A}$.

  (1) $\Phi$ contains $\Zero$.
  \begin{proof}
    $\Zero \NotElem A$.
    Hence $\Zero$ is less than every element of $A$.
    Thus $\Zero \Elem \Phi$.
  \end{proof}

  (2) For all $n \Elem \Phi$ we have $n \Plus \One \Elem \Phi$.
  \begin{proof}
    Let $n \Elem \Phi$.
    Then $n$ is less than any element of $A$.
    Assume that $\Phi$ does not contain $n \Plus \One$.
    Then we can take an $m \Elem A$ such that $n \Plus \One \NotLess m$ (by a).
    Then $n \Less m \Leq n \Plus \One$.
    Hence $m \Eq n \Plus \One$.
    Contradiction.
  \end{proof}

  Then $\Phi$ contains every natural number (by \sr{induction I}{induction}).
  Therefore every natural number is less than any element of $A$.
  Consequently $A$ has no elements.
  Contradiction.
\end{proof}

\begin{theorem*}[forthel,name=induction II]
  Let $A$ be a class.
  Assume for all $n \Elem \Naturals$ if $A$ contains all predecessors of $n$ then $A$ contains $n$.
  Then $A$ contains every natural number.
\end{theorem*}
\begin{proof}[forthel]
  Assume the contrary.
  Take a natural number $n$ that is not contained in $A$.
  Then $n$ is contained in $\Naturals \Diff A$.
  Hence we can take a $m \Elem \Naturals \Diff A$ such that $m \Leq k$ for all $k \Elem \Naturals \Diff A$.
  Then $\Naturals \Diff A$ does not contain any predecessor of $m$.
  Therefore $A$ contains all predecessors of $m$.
  Consequently $A$ contains $m$.
  Contradiction.
\end{proof}

\begin{theorem*}[forthel,name=induction III]
  Let $A$ be a class.
  Let $k$ be a natural number such that $k \Elem A$.
  Assume that for all $n \Elem \NatGeq{k}$ if $n \Elem A$ then $n \Plus \One \Elem A$.
  Then for all $n \Elem \NatGeq{k}$ we have $n \Elem A$.
\end{theorem*}
\begin{proof}[forthel]
  Define $\Phi \DefEq \SClass{n}{\Naturals}{\text{ if }n \Geq k\text{ then }n \Elem A}$.

  (1) $\Phi$ contains $\Zero$.
  Indeed if $\Zero \Geq k$ then $\Zero \Eq k \Elem A$.

  (2) For all $n \Elem \Phi$ we have $n \Plus \One \Elem \Phi$.
  \begin{proof}
    Let $n \Elem \Phi$.

    Let us show that if $n \Plus \One \Geq k$ then $n \Plus \One \Elem A$.
      Assume $n \Plus \One \Geq k$.

      \begin{case}{$n \Less k$.}
        Then $n \Plus \One \Eq k$.
        Hence $n \Plus \One \Elem A$.
      \end{case}

      \begin{case}{$n \Geq k$.}
        Then $n \Elem A$.
        Hence $n \Plus \One \Elem A$.
      \end{case}
    End.

    Therefore $n \Plus \One \Elem \Phi$.
  \end{proof}

  Thus $\Phi$ contains every natural number (by \sr{induction I}{induction}).
  Consequently for all $n \Elem \NatGeq{k}$ we have $n \Elem A$.
\end{proof}

\end{smodule}
\end{document}
